\documentclass[11pt]{article}

\usepackage[margin=1.1in]{geometry}
\usepackage[T1]{fontenc}
\usepackage{amsmath,amssymb,amsthm}
\usepackage{bm}
\usepackage{booktabs}
\usepackage{algorithm}
\usepackage{algpseudocode}
\usepackage[colorlinks=true,linkcolor=blue,citecolor=blue,urlcolor=blue]{hyperref}

% ---------------------------------------------------------------------------
% Notation macros
% ---------------------------------------------------------------------------
\newcommand{\w}{\mathbf{w}}
\newcommand{\vv}{\mathbf{v}}
\newcommand{\XM}{X_{\mathrm{M}}}
\newcommand{\fM}{\mathbf{f}_{\mathrm{M}}}
\newcommand{\mM}{\mathbf{m}_{\mathrm{M}}}
\newcommand{\KMM}{\mathbf{K}_{\mathrm{MM}}}
\newcommand{\gp}{\mathcal{GP}}
\newcommand{\D}{\mathcal{D}}
\newcommand{\N}{\mathcal{N}}
\newcommand{\gbar}{\bar{\mathbf{g}}}
\newcommand{\vhat}{\hat{\mathbf{v}}}
\newcommand{\ghat}{\hat{\mathbf{g}}}
\newcommand{\btau}{\bm{\tau}}
\newcommand{\gradU}{\nabla \widetilde{U}}

\title{A Scale-Adapted, Cyclical Functional SGHMC Sampler:\\
Implementation Notes and Provenance of Modifications}
\author{}
\date{}

\begin{document}
\maketitle

\begin{abstract}
This note documents the functional stochastic gradient Hamiltonian Monte Carlo
(fSGHMC) sampler implemented in this codebase.  The algorithm follows the
functional SG-MCMC framework of Wu et al.~\cite{wu2025},
in which a Gaussian-process (GP) prior over \emph{functions} replaces the usual
weight-space Gaussian prior of a Bayesian neural network (BNN), and its
gradient is injected into the parameter-space dynamics through a
vector--Jacobian product at a random set of measurement points.  The
implementation departs from the reference algorithm of \cite{wu2025} in several
ways: the base sampler is the \emph{scale-adapted} (preconditioned) SGHMC of
Springenberg et al.~\cite{springenberg2016} rather than a vanilla
leapfrog SGHMC; the step size follows the cyclical cosine schedule of Zhang et
al.~\cite{zhang2020} with cool-phase sample harvesting and stationary-law
momentum refreshment; the posterior is tempered by a temperature parameter
$T$; the GP prior gradient is obtained from a double-precision Cholesky solve
of the nugget-regularised prior Gram matrix; and a small set of numerical
safeguards
(gradient-norm clipping, preconditioner floors, an optional per-step
displacement clamp) protect the chain during phase transitions.  We describe
each modification, give its origin, and present pseudocode in the style of the
original paper.
\end{abstract}

% ===========================================================================
\section{Setting and notation}
% ===========================================================================

Let $f(\cdot;\w):\mathcal{X}\to\mathbb{R}$ be the function computed by a
neural network with parameters $\w\in\mathbb{R}^k$, and let
$\D=\{(x_i,y_i)\}_{i=1}^{N}$ be a dataset with likelihood
$p(y\,|\,f(x;\w))$.\footnote{The implementation is agnostic to the form of the
likelihood; any differentiable $\log p(y\,|\,f)$ that depends on the input
only through the network output fits the scheme unchanged.}
Instead of a weight-space prior $p(\w)$, a functional prior
$P_0 \sim \gp(m,K)$ is placed directly on $f$, and the target is the
posterior measure over functions
\begin{equation}
  P_{f|\D}(\mathrm{d}f) \;\propto\;
  \exp\!\bigl(-\Phi(f)\bigr)\,P_0(\mathrm{d}f),
  \qquad
  \Phi(f) = -\log p(Y_\D \,|\, f(X_\D;\w)),
\end{equation}
following the function-space formulation of \cite{wu2025}.  Since $f$ is
determined by $\w$, sampling proceeds in parameter space with the potential
\begin{equation}
  U(\w) \;=\; -\log p(Y_\D\,|\,f(X_\D;\w)) \;+\; I_0\bigl(f(\cdot;\w)\bigr),
  \label{eq:potential}
\end{equation}
where $I_0$ is the Onsager--Machlup functional of $P_0$ (heuristically,
$-\log P_0$).  The prior term is made tractable by evaluating it on a finite
set of \emph{measurement points}
$\XM = [x_1^{\mathrm{m}},\dots,x_M^{\mathrm{m}}]^\top$, freshly drawn at every
step, where the GP marginal is a multivariate Gaussian.  Writing
$\fM = f(\XM;\w)$, $\mM = m(\XM)$, $\KMM = K(\XM,\XM)$, the stochastic
estimate of $\nabla_\w U$ used at each step is
\begin{equation}
  \gradU(\w)
  \;=\;
  \underbrace{\frac{N}{n}\sum_{i\in\mathcal{B}}
      \nabla_\w \bigl[-\log p(y_i\,|\,f(x_i;\w))\bigr]}_{\text{minibatch likelihood
      gradient},\ |\mathcal{B}| = n}
  \;+\;
  \underbrace{J_\w(\XM)^\top\,
      \bigl(\KMM + \tilde\delta I\bigr)^{-1}\bigl(\fM - \mM\bigr)}_{
      \;-\,\nabla_\w \log P_0(\fM)\;=\;J^\top\bm{\alpha}},
  \label{eq:gradU}
\end{equation}
where $J_\w(\XM)\in\mathbb{R}^{M\times k}$ is the Jacobian of the network
outputs at the measurement points with respect to $\w$ and $\tilde\delta$ is a
diagonal regulariser (Section~\ref{sec:priorsolve}).  This is exactly the drift
of fSGHMC in Eq.~(9) of \cite{wu2025}; only the likelihood term is
minibatched, and the prior-gradient term enters unscaled, at $O(1)$ relative
to the full-data likelihood gradient.  The Jacobian is never formed
explicitly: $J^\top\bm{\alpha}$ is one reverse-mode vector--Jacobian product
(a single backward pass with $\bm{\alpha}$ as the output cotangent), with
$\bm{\alpha}$ computed numerically and detached from the autodiff graph.

% ===========================================================================
\section{The implemented algorithm}
% ===========================================================================

\subsection{Base sampler: scale-adapted SGHMC}
\label{sec:adaptive}

The reference algorithm of \cite{wu2025} (their Algorithm~2) is a vanilla
SGHMC in the sense of Chen et al.~\cite{chen2014}: identity mass matrix,
hand-tuned step size and friction, and an inner leapfrog loop of $m$ steps
with the momentum resampled from $\N(0,M)$ at every outer iteration.  The
implementation replaces this with the \emph{scale-adapted} SGHMC of
Springenberg et al.~\cite{springenberg2016} (the sampler underlying
BOHAMIANN), which removes the sensitivity to the raw gradient scale by
adapting a diagonal preconditioner during burn-in.

The sampler maintains, per parameter element and with all operations
element-wise, an exponential moving average $\ghat \approx \mathbb{E}[\gradU]$
of the gradient, an uncentred second-moment estimate
$\vhat \approx \mathbb{E}[(\gradU)^2]$, and an automatic averaging window
$\btau$ whose update rule originates in the adaptive learning-rate method of
Schaul et al.~\cite{schaul2013}.  During the $N_b$ burn-in steps, for
stochastic gradient $\gbar$:
\begin{align}
  \btau &\leftarrow
    \Bigl(1 - \frac{\ghat^{2}}{\vhat + \lambda}\Bigr)\,\btau + 1,
    \qquad \btau \leftarrow \max(\btau,\,1),
    \label{eq:tau}\\
  \ghat &\leftarrow \ghat + \btau^{-1}(\gbar - \ghat),\\
  \vhat &\leftarrow \vhat + \btau^{-1}(\gbar^{2} - \vhat),
    \qquad \vhat \leftarrow \max(\vhat,\,v_{\min}),
    \label{eq:vhat}
\end{align}
with a small numerical stabiliser $\lambda$.  After burn-in the estimates are
frozen, so the recorded samples are drawn from unchanged (valid) dynamics, as
argued in \cite{springenberg2016}.  The clamps in \eqref{eq:tau} and
\eqref{eq:vhat} are safeguards added in this codebase; see
Section~\ref{sec:safeguards}.

Every step then applies the preconditioned update of
\cite{springenberg2016} (their Eq.~10), written in the incremental variable
$\vv = \varepsilon M^{-1}\mathbf{r}$ (momentum after the variable
substitution), with preconditioner
$\mathbf{a} = \bigl(\sqrt{\vhat} + \lambda\bigr)^{-1} \approx
\hat V^{-1/2}$, step size $\varepsilon$, and momentum decay
$c$ (the implementation's \texttt{mdecay}, playing the role of
$\varepsilon \hat V^{-1/2} C$):
\begin{align}
  \bm{\sigma}^{2} &= \max\bigl(2\,\varepsilon^{2} c\,\mathbf{a}
      \;-\; \varepsilon^{4},\ 10^{-16}\bigr),
      \label{eq:noisevar}\\
  \vv &\leftarrow \vv \;-\; \varepsilon^{2}\,\mathbf{a}\odot\gbar
      \;-\; c\,\vv \;+\; \bm{\sigma}\odot\bm{\xi},
      \qquad \bm{\xi}\sim\N(0,I),
      \label{eq:momentum}\\
  \w &\leftarrow \w + \vv.
  \label{eq:position}
\end{align}
The $-\varepsilon^4$ term in \eqref{eq:noisevar} is the correction for the
minibatch-noise estimate $\hat B = \tfrac{1}{2}\varepsilon\hat V$, which
cancels against the squared preconditioner exactly as in
\cite{springenberg2016}.

Two structural differences from Algorithm~2 of \cite{wu2025} follow from this
choice of base sampler:
\begin{enumerate}
  \item \textbf{No inner leapfrog loop.}  The chain is a single persistent
  trajectory in the style of \cite{chen2014}: one
  momentum/parameter update per gradient evaluation, with friction $c$
  providing the mixing that periodic momentum resampling provides in
  \cite{wu2025}.  Momentum is refreshed only at cycle boundaries
  (Section~\ref{sec:cyclical}).
  \item \textbf{Diagonal preconditioning.}  The mass matrix is the adapted
  $M^{-1} = \mathrm{diag}(\hat V^{-1/2})$ rather than the identity, making
  the effective step size per parameter scale-invariant.
\end{enumerate}

\subsection{Tempering}
\label{sec:tempering}

The gradient actually passed to updates
\eqref{eq:tau}--\eqref{eq:momentum} is
\begin{equation}
  \gbar \;=\; \frac{1}{T}\,\gradU(\w),
\end{equation}
where $T$ is a posterior temperature (implemented as the sampler's gradient
scale $N/T$ applied to a per-example-normalised loss).  Since the injected
noise in \eqref{eq:noisevar} is unchanged, the chain targets the tempered
posterior $\propto \bigl[p(Y_\D|f)\,P_0\bigr]^{1/T}$; $T=1$ recovers the
untempered target.  Tempering the posterior of a deep BNN is common practice
in SG-MCMC --- Zhang et al.~\cite{zhang2020} introduce a system temperature
for the same sampler family, and Wenzel et al.~\cite{wenzel2020} study the
effect systematically ("cold posteriors", $T<1$).  Note that both the
likelihood and prior terms are tempered jointly.

\subsection{Cyclical step-size schedule with cool-phase harvesting}
\label{sec:cyclical}

After burn-in the step size follows the cyclical cosine schedule of Zhang et
al.~\cite{zhang2020} (itself adapted from warm-restart schedules in
optimisation \cite{loshchilov2017}).  With cycle length $L$, post-burn-in
step index $j = (t - N_b) \bmod L$, and bounds
$\varepsilon_{\min},\varepsilon_{\max}$:
\begin{equation}
  \varepsilon_t \;=\; \varepsilon_{\min}
  + \tfrac{1}{2}\bigl(\varepsilon_{\max} - \varepsilon_{\min}\bigr)
    \Bigl(1 + \cos\tfrac{\pi j}{L}\Bigr).
  \label{eq:cyclical}
\end{equation}
Each cycle begins with large, exploratory steps and anneals toward
$\varepsilon_{\min}$; samples are recorded only near the cold end.  The
implementation differs from \cite{zhang2020} in three deliberate ways:
\begin{enumerate}
  \item \textbf{Nonzero floor.}  The schedule anneals to
  $\varepsilon_{\min} > 0$ rather than to $0$, so the chain keeps sampling
  (with small steps) instead of degenerating to optimisation; relatedly,
  there is no explicit $T=0$ "optimisation" stage inside the exploration
  phase of a cycle --- exploration is realised purely through the large step
  size.  A single conventional burn-in phase precedes the first cycle (during
  which the preconditioner of Section~\ref{sec:adaptive} adapts), rather than
  treating each cycle's exploration stage as its own burn-in.
  \item \textbf{Cool-phase harvesting.}  Let $\rho\in(0,1]$ be the cool
  fraction and $S_c \ge 1$ the samples harvested per cycle.  With cool-window
  length $\ell = \max(1, \lfloor \rho L\rfloor)$ and thinning interval
  $\Delta = \max(1, \lfloor \ell / S_c\rfloor)$, a sample is recorded at step
  $j$ iff the distance from the cycle end, $d = (L-1) - j$, satisfies
  $d \in \{0, \Delta, 2\Delta, \dots, (S_c - 1)\Delta\}$ --- i.e., $S_c$
  thinned samples spaced to terminate exactly at the coldest step of the
  cycle.  $S_c = 1$ keeps only the coldest iterate per cycle.
  \item \textbf{Stationary-law momentum refreshment.}  At each cycle start
  ($j=0$) the momentum variable is resampled from its stationary distribution
  under the preconditioned dynamics,
  \begin{equation}
    \vv \;\sim\; \N\!\bigl(0,\ \varepsilon_t^{2}\,
      \mathrm{diag}(\mathbf{a})\bigr),
    \qquad
    \mathbf{a} = \bigl(\sqrt{\vhat}+\lambda\bigr)^{-1},
    \label{eq:refresh}
  \end{equation}
  with $\varepsilon_t = \varepsilon_{\max}$ at the cycle start (for
  unpreconditioned SGHMC this reduces to
  $\vv \sim \N(0, \varepsilon_t I)$).  This is the incremental-variable
  analogue of the $z_t \sim \N(0,M)$ refreshment in Algorithm~2 of
  \cite{wu2025}: since $\vv = \varepsilon M^{-1} \mathbf{z}$ with
  $M^{-1} = \mathrm{diag}(\mathbf a)$ and $\mathbf z\sim\N(0,M)$, one gets
  $\operatorname{Var}(\vv) = \varepsilon^2 M^{-1}$.  Resampling (rather than
  zeroing) the momentum avoids dropping the chain into an atypical set that
  would cost $O(1/c)$ steps to re-equilibrate.
\end{enumerate}
When the cyclical schedule is disabled, the sampler falls back to a constant
step size with fixed-interval thinning (record every $\kappa$-th post-burn-in
iterate), matching the "draw separated samples every $\kappa$ epochs"
protocol of \cite{wu2025}.

\subsection{The GP prior solve}
\label{sec:priorsolve}

The prior-gradient cotangent
$\bm{\alpha} = (\KMM + \tilde\delta I)^{-1}(\fM - \mM)$ in
Eq.~\eqref{eq:gradU} is computed by a direct Cholesky factorisation of the
$M \times M$ Gram matrix, carried out in double precision and detached from
the autodiff graph (only the subsequent VJP touches the network's
computation graph).  The prior kernels used in practice are full rank by
construction: the kernel includes an explicit nugget (observation-noise)
term $\sigma_n^2\,\delta_{x,x'}$ on the diagonal, corresponding to the valid
GP kernel $k_{\mathrm{eff}}(x,x') = K(x,x') + \sigma_n^2\,\delta_{x,x'}$
\cite{rasmussen2006}, so the factorisation is well conditioned without
further regularisation ($\tilde\delta$ in Eq.~\eqref{eq:gradU} absorbs the
nugget plus any optional extra jitter).  One convention is worth noting: a
global prior-amplitude hyperparameter multiplies the \emph{entire} Gram
matrix, nugget included ($\KMM \to a^2 \KMM$), so the conditioning of the
Cholesky solve is invariant to the amplitude and the prior-gradient
magnitude scales exactly as $1/a^{2}$.  If the measurement set is empty the
prior term is skipped entirely and the update reduces to likelihood-only
SGHMC.

\subsection{Numerical safeguards}
\label{sec:safeguards}

Four safeguards, none of which appear in
\cite{wu2025,springenberg2016,zhang2020}, protect the chain --- chiefly at
the transition from burn-in to the first hot phase of the cyclical schedule,
where the preconditioner calibrated on burn-in gradients suddenly meets
$\varepsilon_{\max}$-sized steps:

\begin{itemize}
  \item \textbf{Second-moment floor} ($v_{\min}$, Eq.~\ref{eq:vhat}).  A
  parameter that receives near-zero gradient during burn-in (e.g.\ a dead
  ReLU unit) drives $\vhat\to 0$, and with only the tiny stabiliser
  $\lambda$ in the denominator the preconditioner gain
  $1/(\sqrt{\vhat}+\lambda)$ can reach $\sim\!10^{8}$, producing a
  catastrophic jump on the first nonzero gradient.  Flooring
  $\vhat \ge v_{\min}$ bounds the gain by $v_{\min}^{-1/2}$.
  \item \textbf{Averaging-window clamp} ($\btau \ge 1$,
  Eq.~\ref{eq:tau}).  Floating-point rounding can push $\btau$ below $1$
  when gradients are small; a negative $\btau^{-1}$ would reverse the
  moving-average updates and can drive $\vhat$ negative
  ($\sqrt{\vhat} = \mathrm{NaN}$).
  \item \textbf{Gradient-norm clipping.}  The total gradient norm (of the
  un-tempered, per-example-scale gradient) is clipped to a threshold
  $\kappa_{\mathrm{clip}}$ --- always during burn-in, and during the sampling
  phase only if explicitly enabled, since clipping while samples are being
  recorded biases the sampled posterior (in particular its tails).  The
  pre-clip norm is instrumented in both phases so that runs can verify the
  clip never fires while sampling.
  \item \textbf{Displacement clamp} (optional).  Each element of $\vv$ may be
  clamped to $[-v_{\mathrm{max}}, v_{\mathrm{max}}]$ before the position
  update \eqref{eq:position}, bounding the per-step parameter change.  The
  steady-state momentum scale is $O(\varepsilon^{2}/c)$; the clamp is set
  well above that level so it fires only in pathological transients.
\end{itemize}

\subsection{Initialisation and parallel chains}

Rather than initialising $\w_0 \sim \N(0,I)$ as in \cite{wu2025}, chains
start from a shared warm-up point (e.g.\ the end state of a preliminary
burn-in-only run).  When $R$ chains are run in parallel, chain $r$ perturbs
every parameter tensor of the shared start by
$\N\bigl(0, (\iota\cdot\mathrm{std}(\w))^{2}\bigr)$ with a relative jitter
scale $\iota$, giving overdispersed initial states so that pooled-chain
convergence diagnostics such as $\widehat R$ remain valid and the chains can
cover distinct posterior modes~\cite{gelman1992}.  Chains are otherwise fully
independent (separate seeds, RNG streams, and sample stores).

% ===========================================================================
\section{Summary of modifications and their origins}
% ===========================================================================

\begin{table}[h]
\centering
\small
\begin{tabular}{p{0.46\textwidth} p{0.46\textwidth}}
\toprule
\textbf{Component / modification} & \textbf{Origin} \\
\midrule
Functional prior gradient at random measurement points; parameter-space
functional dynamics & Wu et al.~\cite{wu2025} (fSGHMC) \\
\addlinespace
Persistent-momentum SGHMC with friction (no leapfrog inner loop, no MH
correction) & Chen et al.~\cite{chen2014} \\
\addlinespace
Diagonal preconditioner $M^{-1} = \mathrm{diag}(\hat V^{-1/2})$ adapted
during burn-in; noise-estimate cancellation; frozen after burn-in &
Springenberg et al.~\cite{springenberg2016} (BOHAMIANN) \\
\addlinespace
Automatic averaging window $\btau$ for the moment estimates &
Schaul et al.~\cite{schaul2013}, via \cite{springenberg2016} \\
\addlinespace
Cyclical cosine step size; exploration/sampling stages; samples from the cold
phase & Zhang et al.~\cite{zhang2020} (cSG-MCMC); cosine restarts from
\cite{loshchilov2017} \\
\addlinespace
Posterior temperature $T$ & Zhang et al.~\cite{zhang2020}; see also Wenzel et
al.~\cite{wenzel2020} \\
\addlinespace
Momentum refreshment at cycle starts from the preconditioned stationary law
\eqref{eq:refresh} & Refreshment per Algorithm~2 of \cite{wu2025}; the
preconditioned form and cycle-start placement are this codebase's \\
\addlinespace
Multi-sample cool-phase harvesting ($S_c$ thinned samples ending at the
coldest step) & This codebase \\
\addlinespace
Double-precision Cholesky solve of the nugget-regularised Gram;
amplitude-invariant conditioning convention & Standard GP practice
\cite{rasmussen2006}; conventions are this codebase's \\
\addlinespace
Safeguards: $v_{\min}$ floor, $\btau\ge1$ clamp, phase-scoped gradient
clipping, displacement clamp & This codebase (heuristics; not from the
literature) \\
\addlinespace
Warm-started, overdispersed parallel chains & Standard MCMC practice
\cite{gelman1992} \\
\bottomrule
\end{tabular}
\caption{Provenance of each component of the implemented sampler.}
\label{tab:origins}
\end{table}

% ===========================================================================
\section{Pseudocode}
% ===========================================================================

Algorithm~\ref{alg:step} gives the per-iteration update (the scale-adapted
SGHMC step applied to the functional-posterior gradient), and
Algorithm~\ref{alg:fsghmc} the full sampler, in the style of Algorithm~2 of
\cite{wu2025}.  All vector operations are element-wise.

\begin{algorithm}[t]
\caption{\textsc{AdaptiveStep}: scale-adapted SGHMC update
\cite{springenberg2016}}
\label{alg:step}
\begin{algorithmic}[1]
\Require gradient $\gbar$; step size $\varepsilon$; momentum decay $c$;
  state $(\btau, \ghat, \vhat, \vv)$; step counter $t$; burn-in length $N_b$;
  stabiliser $\lambda$; floor $v_{\min}$; optional clamp $v_{\mathrm{max}}$
\If{$t \le N_b$}  \Comment{burn-in: adapt preconditioner}
  \State $\btau \gets \bigl(1 - \ghat^{2}/(\vhat+\lambda)\bigr)\btau + 1$;
    \quad $\btau \gets \max(\btau, 1)$
  \State $\ghat \gets \ghat + \btau^{-1}(\gbar - \ghat)$
  \State $\vhat \gets \vhat + \btau^{-1}(\gbar^{2} - \vhat)$;
    \quad $\vhat \gets \max(\vhat, v_{\min})$
\EndIf
\State $\mathbf{a} \gets \bigl(\sqrt{\vhat} + \lambda\bigr)^{-1}$
  \Comment{$\approx \hat V^{-1/2}$}
\State $\bm{\sigma}^{2} \gets \max\bigl(2\varepsilon^{2} c\,\mathbf{a}
  - \varepsilon^{4},\, 10^{-16}\bigr)$
\State $\vv \gets \vv - \varepsilon^{2}\mathbf{a}\odot\gbar - c\,\vv
  + \bm{\sigma}\odot\bm{\xi}$, \quad $\bm{\xi} \sim \N(0, I)$
\If{$v_{\mathrm{max}}$ set}
  $\vv \gets \mathrm{clamp}(\vv, -v_{\mathrm{max}}, v_{\mathrm{max}})$
\EndIf
\State $\w \gets \w + \vv$
\end{algorithmic}
\end{algorithm}

\begin{algorithm}[t]
\caption{Cyclical scale-adapted functional SGHMC (this codebase)}
\label{alg:fsghmc}
\begin{algorithmic}[1]
\Require dataset $\D$ ($|\D| = N$), minibatch size $n$; GP prior
  $\gp(m, K)$ with jitter $\delta$; measurement pool
  $\mathcal{X}_{\mathrm{pool}}$ and draw size $M$; temperature $T$;
  burn-in steps $N_b$; cycle length $L$; step sizes
  $\varepsilon_{\min}, \varepsilon_{\max}$; cool fraction $\rho$;
  samples per cycle $S_c$; discard count $D_0$; target sample count
  $S_{\mathrm{tot}}$; momentum decay $c$; clip threshold
  $\kappa_{\mathrm{clip}}$ (burn-in only, unless enabled for sampling);
  warm start $\w^{\star}$, jitter scale $\iota$
\State $\w_0 \gets \w^{\star} + \iota\,\mathrm{std}(\w^{\star})\odot
  \bm{\zeta}$, \quad $\bm{\zeta}\sim\N(0,I)$
  \Comment{overdispersed init (per chain)}
\State $(\btau, \ghat, \vhat) \gets \mathbf{1}$;\quad
  $\vv \gets \mathbf{0}$;\quad
  $\mathcal{S} \gets \emptyset$;\quad
  $s \gets 0$
\For{$t = 0, 1, 2, \dots$ until $|\mathcal{S}| = S_{\mathrm{tot}}$}
  \State \textbf{\# Step size and momentum refreshment}
  \If{$t < N_b$} \quad $\varepsilon_t \gets \varepsilon_{\min}$
  \Else
    \State $j \gets (t - N_b) \bmod L$;\qquad
      $\varepsilon_t \gets \varepsilon_{\min} + \tfrac12
      (\varepsilon_{\max} - \varepsilon_{\min})
      \bigl(1 + \cos\tfrac{\pi j}{L}\bigr)$
    \If{$j = 0$} \quad
      $\vv \sim \N\bigl(0,\ \varepsilon_t^{2}\,
        \mathrm{diag}\bigl((\sqrt{\vhat}+\lambda)^{-1}\bigr)\bigr)$
      \Comment{Eq.~\eqref{eq:refresh}}
    \EndIf
  \EndIf
  \State \textbf{\# Stochastic functional-posterior gradient}
  \State draw minibatch $\mathcal{B} \subset \D$, $|\mathcal{B}| = n$;\qquad
    $\mathbf{g}_{\mathrm{lik}} \gets \tfrac{1}{n} \sum_{i\in\mathcal{B}}
    \nabla_\w\bigl[-\log p(y_i\,|\,f(x_i;\w_t))\bigr]$
  \State draw measurement set $\XM \subset \mathcal{X}_{\mathrm{pool}}$
    without replacement, $|\XM| = M$
  \State $\fM \gets f(\XM;\w_t)$;\qquad
    $\bm{\alpha} \gets (\KMM + \tilde\delta I)^{-1}(\fM - \mM)$
    \Comment{Cholesky, \S\ref{sec:priorsolve}}
  \State $\mathbf{g}_{\mathrm{prior}} \gets J_{\w}(\XM)^{\top}\bm{\alpha}$
    \Comment{one VJP backward pass}
  \State $\gbar \gets \mathbf{g}_{\mathrm{lik}}
    + \mathbf{g}_{\mathrm{prior}} / N$
  \If{$t < N_b$ \textbf{or} clip-in-sampling enabled}
    \State $\gbar \gets \gbar \cdot
      \min\bigl(1,\ \kappa_{\mathrm{clip}} / \lVert\gbar\rVert_2\bigr)$
  \EndIf
  \State $\gbar \gets (N / T)\,\gbar$
    \Comment{$= \tfrac{1}{T}\gradU(\w_t)$, Eq.~\eqref{eq:gradU}}
  \State \textbf{\# Update}
  \State $\w_{t+1} \gets$ \textsc{AdaptiveStep}$(\gbar, \varepsilon_t, c,
    \text{state}, t, N_b)$ \Comment{Algorithm~\ref{alg:step}}
  \State \textbf{\# Cool-phase sample harvesting}
  \If{$t \ge N_b$}
    \State $\ell \gets \max(1, \lfloor\rho L\rfloor)$;\quad
      $\Delta \gets \max(1, \lfloor\ell / S_c\rfloor)$;\quad
      $d \gets (L - 1) - j$
    \If{$d \in \{0, \Delta, \dots, (S_c{-}1)\Delta\}$}
      \State $s \gets s + 1$;\qquad
        \textbf{if} $s > D_0$ \textbf{then}
        $\mathcal{S} \gets \mathcal{S} \cup \{\w_{t+1}\}$
    \EndIf
  \EndIf
\EndFor
\State \Return $\mathcal{S}$
\end{algorithmic}
\end{algorithm}

As in \cite{wu2025}, posterior predictions integrate over the collected
(pooled, over chains) samples:
$p(y^{*}|x^{*},\D) \approx \tfrac{1}{|\mathcal{S}|}
\sum_{\w\in\mathcal{S}} p\bigl(y^{*}\,|\,f(x^{*};\w)\bigr)$.

% ===========================================================================
\begin{thebibliography}{10}

\bibitem{wu2025}
M.~Wu, J.~Xuan, and J.~Lu.
\newblock Functional stochastic gradient {MCMC} for {B}ayesian neural
  networks.
\newblock In \emph{Proceedings of the 28th International Conference on
  Artificial Intelligence and Statistics (AISTATS)}, PMLR 258, 2025.

\bibitem{chen2014}
T.~Chen, E.~Fox, and C.~Guestrin.
\newblock Stochastic gradient {H}amiltonian {M}onte {C}arlo.
\newblock In \emph{Proceedings of the 31st International Conference on
  Machine Learning (ICML)}, 2014.

\bibitem{springenberg2016}
J.~T. Springenberg, A.~Klein, S.~Falkner, and F.~Hutter.
\newblock Bayesian optimization with robust {B}ayesian neural networks.
\newblock In \emph{Advances in Neural Information Processing Systems 29
  (NeurIPS)}, 2016.

\bibitem{schaul2013}
T.~Schaul, S.~Zhang, and Y.~LeCun.
\newblock No more pesky learning rates.
\newblock In \emph{Proceedings of the 30th International Conference on
  Machine Learning (ICML)}, 2013.

\bibitem{zhang2020}
R.~Zhang, C.~Li, J.~Zhang, C.~Chen, and A.~G. Wilson.
\newblock Cyclical stochastic gradient {MCMC} for {B}ayesian deep learning.
\newblock In \emph{International Conference on Learning Representations
  (ICLR)}, 2020.

\bibitem{loshchilov2017}
I.~Loshchilov and F.~Hutter.
\newblock {SGDR}: Stochastic gradient descent with warm restarts.
\newblock In \emph{International Conference on Learning Representations
  (ICLR)}, 2017.

\bibitem{wenzel2020}
F.~Wenzel, K.~Roth, B.~S. Veeling, J.~{\'S}wi{\k{a}}tkowski, L.~Tran,
  S.~Mandt, J.~Snoek, T.~Salimans, R.~Jenatton, and S.~Nowozin.
\newblock How good is the {B}ayes posterior in deep neural networks really?
\newblock In \emph{Proceedings of the 37th International Conference on
  Machine Learning (ICML)}, 2020.

\bibitem{rasmussen2006}
C.~E. Rasmussen and C.~K.~I. Williams.
\newblock \emph{Gaussian Processes for Machine Learning}.
\newblock MIT Press, 2006.

\bibitem{gelman1992}
A.~Gelman and D.~B. Rubin.
\newblock Inference from iterative simulation using multiple sequences.
\newblock \emph{Statistical Science}, 7(4):457--472, 1992.

\end{thebibliography}

\end{document}
